\chapter{Diskussion und Ausblick}
\label{ch:DiskussionAusblick}


\section{Interpretation der Ergebnisse}
\label{sec:InterpretationErgebnisse}
\todo[inline]{Interpretation der Ergebnisse}


\section{Limitationen und Herausforderungen}
\label{sec:LimitationenHerausforderungen}
\todo[inline]{Limitationen und Herausforderungen}


\section{Weiterführende Konzepte}
\label{sec:Konzepte}
Im Zuge des Projekts wurden drei verschiedene Konzepte erarbeitet, die darauf abzielen, den Generator durch 
zusätzliche Trainingsstrategien zu verbessern und dadurch die Qualität der erzeugten Daten für nachgelagerte 
Verfahren zu erhöhen. Die folgenden Abschnitte stellen diese Konzepte mit ihren jeweiligen Grundideen sowie den 
identifizierten Vor- und Nachteilen vor. Nach eingehender Diskussion würde die Projektgruppe, sollte das Thema in 
zukünftigen Arbeiten weiterverfolgt werden, insbesondere Konzept~3 als vielversprechendste Option betrachten.

\subsection{Konzept 1: Feature-Anzahl}
\label{sec:Konzept1}
Das erste Kopplungskonzept zwischen Generator und \ac{xFeat} basiert auf der Annahme, dass eine höhere Anzahl erkannter Merkmalspunkte eine bessere Bildrepräsentation anzeigt. 
Der Generator erzeugt dazu augmentierte Bilddaten, die von \ac{xFeat} verarbeitet werden. 
Aus den entstehenden Keypoint-Heatmaps wird über eine soft-differenzierbare Approximation (z.~B. Sigmoid-Gewichtung oberhalb eines adaptiven Schwellwerts) ein „soft count“ berechnet. 
Dieser Wert wird mit einer Zielanzahl von Merkmalen verglichen und als zusätzlicher Loss-Term in das Training integriert. 
Der Generator wird belohnt, wenn genügend Features detektiert werden, und bestraft, wenn die Anzahl zu gering ist.  

Ähnliche Ansätze existieren bereits: Arbeiten wie UKPGAN~\cite{you2022ukpgan} steuern die Verteilung von Keypoints über generative Modelle, während RIPE~\cite{kunzel2025ripe} Reward-Signale ohne Ground Truth nutzt, um Anzahl und Qualität zu regulieren.  

Für dieses Projekt – die zuverlässige Bestimmung von Greifpositionen bei texturarmen Bauteilen – zeigt das Konzept Vor- und Nachteile. 
Ein Vorteil ist die einfache Integration in den Trainingsloop und die Nutzung von \ac{xFeat} ohne Änderungen. 
Kritisch ist jedoch, dass die Keypoint-Anzahl allein kein Qualitätskriterium ist: 
Der Generator könnte künstliche Texturen erzeugen („Reward Hacking“), was dem Ziel widerspricht. 
Daher sind regulierende Zusatzverluste notwendig.  

Insgesamt eignet sich Konzept~1 als Machbarkeitsnachweis einer differentiellen Kopplung zwischen Generator und Detektor. 
Für die gewünschte Fokussierung auf kantenbasierte Merkmale ist es jedoch unzureichend und wird nicht weiterverfolgt.


\subsection{Konzept 2: Loss}
\label{sec:Konzept2}
Die Grundidee von Konzept~2 besteht darin, Informationen aus der Keypoint-Detektion und -Beschreibung zu nutzen, um den Generator zu steuern. 
Ein niedriger \ac{xFeat}-bezogener Loss gilt als Belohnung, ein hoher Loss als Bestrafung. 
Ziel ist, dass der Generator deformierte Bilder so erzeugt, dass darin robuste und konsistente Features für Matching und Visual Servoing erhalten bleiben~\cite{potje2024xfeat}.  

Für die Umsetzung wird der Generator eng mit \ac{xFeat} gekoppelt: 
Die Generatorausgabe dient als Eingabe für \ac{xFeat}, dessen Keypoints oder Feature-Karten wiederum den Loss bestimmen. 
\ac{xFeat} bleibt eingefroren und liefert ausschließlich Feedback. 
Zwei Varianten sind denkbar:  

\textbf{Black-Box-Variante:} Keypoints und Matches werden ohne Gradientenfluss berechnet. 
Indirekte Größen wie Kantenstärke, Gradientenrichtung oder Patch-Ähnlichkeiten erzeugen einen Score (0–1), der die Konsistenz bewertet. 
Vorteil: robust und einfach; Nachteil: nur indirekte Qualitätsmaße~\cite{szeliski2022}.  

\textbf{White-Box-Variante:} Dichte Karten (z.~B. Heatmaps, Deskriptormaps) fließen in den Loss ein. 
Dadurch wird die Konsistenz direkter bewertet. 
Vorteil: gezieltere Steuerung; Nachteil: rechenintensiv und erfordert Änderungen im Modellcode, da einige Operationen (z.~B. Top-K-Selektion) nicht differenzierbar sind~\cite{detone2018superpoint,sarlin2020superglue}.  

Konzept~2 bietet den Vorteil, dass der Generator durch zusätzliche strukturelle Informationen aus einem etablierten Feature-Detektor gesteuert wird, was die Robustheit für nachgelagerte Aufgaben steigern kann. 
Nachteile sind jedoch der hohe Implementierungsaufwand, mögliche Artefakte bei zu starker Gewichtung des Zusatzlosses und die Gefahr von „Shortcuts“, bei denen der Generator den Loss minimiert, ohne qualitativ bessere Bilder zu erzeugen. 
Deshalb empfiehlt sich eine moderate Gewichtung des \ac{xFeat}-Losses sowie eine Kombination mit weiteren Verlustkomponenten wie adversarialen, rekonstruktiven oder perzeptuellen Losses~\cite{goodfellow2014gan,isola2018Pix2Pix}.  

Die Projektgruppe kam daher zum Schluss, dass Konzept~2 zwar interessante Ansätze bietet, aufgrund der praktischen Herausforderungen und potenziellen Fehlanreize aber nicht priorisiert weiterverfolgt werden sollte.

\subsection{Konzept 3: Keypoint-Position}
\label{sec:Konzept3}
Konzept~3 zielt darauf ab, den Generator anhand der Position erkannter Features zu steuern. 
Keypoints auf stabilen, strukturell relevanten Bereichen (v.~a. Kanten) werden belohnt, während Punkte auf homogenen Flächen bestraft werden. 
Grundlage sind die von \ac{xFeat} gelieferten Koordinaten~\cite{potje2024xfeat}, anhand derer für jedes generierte Bild eine Kantenkarte $E(x,y)$ (z.~B. Sobel, Canny, Harris-Corner) berechnet wird~\cite{szeliski2022}. 
Aus der Kantenstärke und -richtung an den Keypoint-Positionen entsteht ein Score, der misst, wie konsistent die Punkte mit Kanten übereinstimmen. 
Ein hoher Score bedeutet robuste Keypoints, ein niedriger Score ungeeignete Lagen. 
Der normierte Score wird als zusätzlicher Loss-Term in das Training integriert~\cite{chaumette2006visual,chaumette2007visual}.  

Die Idee, Kanten als stabile Ankerpunkte zu nutzen, ist etabliert: 
Schon klassische Verfahren wie \ac{SIFT} oder \ac{SURF} verwenden Gradienteninformationen zur robusten Feature-Detektion~\cite{lowe2004distinctive,bay2006surf}. 
Neuere Arbeiten wie Feature Matching Based on Consistent Edge Structures~\cite{mdpi2020edge} oder SiLK~\cite{gleize2023silk} bestätigen ebenfalls die Robustheit kantenbasierter Strukturen.  

Für das Projektziel – Greifpositionen und Ausrichtungen an texturarmen Bauteilen – ist dieser Ansatz besonders attraktiv, da Kanten weniger anfällig für Beleuchtungsänderungen und Oberflächenrauschen sind.  

Herausforderungen ergeben sich durch die Definition, wann ein Punkt „auf einer Kante“ liegt, da Schwellenwerte und Kantenbreiten variieren können. 
Zudem besteht die Gefahr künstlicher Kantenartefakte zur Maximierung des Rewards. 
Gegenmaßnahmen sind Regularisierungen wie Total Variation oder Frequenzbestrafungen, um unnatürliche Strukturen zu verhindern.  

Trotz dieser Einschränkungen bietet Konzept~3 eine vielversprechende Möglichkeit, die Bildqualität im Hinblick auf robuste und praxisrelevante Feature-Erkennung zu verbessern.
